\documentclass[popl.tex]{subfiles}
\begin{document}

\section{Background}

Recall that a CFG, $\mathcal{G} = \langle \Sigma, V, P, S\rangle$, is a quadruple consisting of terminals $(\Sigma)$, nonterminals $(V)$, productions $\big(P\colon V \rightarrow (V \mid \Sigma)^+\big)$, and a start symbol, $(S)$. Every CFG is reducible to so-called \textit{Chomsky Normal Form}~\cite{chomsky1959certain}, $P'\colon V \rightarrow (V^2 \mid \Sigma)$, where every production is either (1) a binary production $w \rightarrow xz$, or (2) a unit production $w \rightarrow t$, where $w, x, z: V$ and $t: \Sigma$. For example:\vspace{-3pt}

\begin{table}[H]
  \begin{tabular}{llll}
    $G = \big\{\;S \rightarrow S\:S \mid (\:S\:) \mid (\:)\;\big\} \Longrightarrow G' = \big\{\;S\rightarrow Q\:R \mid S\:S \mid L\:R,$ & $R \rightarrow\:),$ & $L \rightarrow (,$ & $Q\rightarrow L\:S\;\big\}$
  \end{tabular}
\end{table}\vspace{-8pt}

Likewise, a finite state automaton (FSA) is a quintuple $\mathcal{A} = \langle Q, \Sigma, \delta, q_\alpha, F\rangle$, where $Q$ is a finite set of states, $\Sigma$ is a finite alphabet, $\delta \subseteq Q \times \Sigma \times Q$ is the transition function, $q_\alpha$ is the initial state, and $F \subseteq Q$ are the accepting states. These generally come in two varieties, deterministic and nondeterministic depending on whether or not $\delta$ maps each pair $\langle q, s \rangle$ to a unique $q'$.

There is an equivalent characterization of the regular languages via an inductively defined datatype, which is often more elegant than FSAs to work with. Consider the generalized regular expression (GRE) fragment containing concatenation, conjunction and disjunction:

\begin{definition}[Star-free GRE fragment]\label{thm:gre_def}
  Let \( e: E \) be an expression defined by the grammar:
  \[
    e \rightarrow \varnothing \mid \varepsilon \mid \Sigma \mid e \cdot e \mid e \lor e \mid e \land e
  \]

where $\varepsilon$ is the empty symbol. Semantically, we interpret these expressions as denoting languages:\vspace{-0.8cm}

  \setlength{\columnseprule}{0pt}
  \setlength{\columnsep}{-3cm}
  \begin{multicols}{2}
    \begin{eqnarray*}
      \mathcal{L}(&\hspace{-0.35cm} \varnothing \hspace{-0.35cm}&) = \varnothing \\
      \mathcal{L}(&\hspace{-0.35cm} \varepsilon \hspace{-0.35cm}&) = \{\varepsilon\} \\
      \mathcal{L}(&\hspace{-0.35cm} a           \hspace{-0.35cm}&) = \{a\}
    \end{eqnarray*} \break\vspace{-0.45cm}
    \begin{eqnarray*}
      \mathcal{L}(&\hspace{-0.35cm} x\cdot z \hspace{-0.35cm}&) = \mathcal{L}(x) \circ \mathcal{L}(z)\text{\footnotemark}\\
      \mathcal{L}(&\hspace{-0.35cm} x\vee  z \hspace{-0.35cm}&) = \mathcal{L}(x) \cup  \mathcal{L}(z)\\
      \mathcal{L}(&\hspace{-0.35cm} x\land z \hspace{-0.35cm}&) = \mathcal{L}(x) \cap  \mathcal{L}(z)
    \end{eqnarray*}
  \end{multicols}
  \footnotetext{Where $\mathcal{L}(x)\circ\mathcal{L}(z)$ is defined as $\big\{a \cdot b \mid a \in \mathcal{L}(x) \land b \in \mathcal{L}(z) \big\}$.}
\end{definition}\vspace{-0.2cm}

\noindent Brzozowski~\cite{brzozowski1964derivatives} introduces an operator, $\partial: E \times \Sigma \rightarrow E$, which quotients a language by some prefix,

\begin{definition}[Brzozowski, 1964]
  To compute the quotient \(\partial_a(L) = \{b \mid ab \in L\}\), we:

  \vspace{-0.8cm}
  \begin{multicols}{2}
    \begin{eqnarray*}
      \phantom{--}\partial_a(&\hspace{-0.35cm} \varnothing \hspace{-0.35cm}&) = \varnothing                                           \\
      \phantom{--}\partial_a(&\hspace{-0.35cm} \varepsilon \hspace{-0.35cm}&) = \varnothing                                           \\
      \phantom{--}\partial_a(&\hspace{-0.35cm} b           \hspace{-0.35cm}&) = \begin{cases}\varepsilon &\text{ if } a = b\\ \varnothing &\text{ if } a \neq b \end{cases}\\
      \phantom{--}\partial_a(&\hspace{-0.35cm} x\cdot z    \hspace{-0.35cm}&) = (\partial_a x)\cdot z \vee \delta(x)\cdot\partial_a z \\
      \phantom{--}\partial_a(&\hspace{-0.35cm} x\vee  z    \hspace{-0.35cm}&) =  \partial_a x \vee  \partial_a z                       \\
      \phantom{--}\partial_a(&\hspace{-0.35cm} x\land z    \hspace{-0.35cm}&) =  \partial_a x \land \partial_a z
    \end{eqnarray*} \break\vspace{-0.45cm}
    \begin{eqnarray*}
      \delta(&\hspace{-0.35cm} \varnothing \hspace{-0.35cm}&) = \varnothing                                      \\
      \delta(&\hspace{-0.35cm} \varepsilon \hspace{-0.35cm}&) = \varepsilon                                      \\
      \delta(&\hspace{-0.35cm} a           \hspace{-0.35cm}&) = \varnothing\phantom{\begin{cases}\varepsilon\\\varnothing\end{cases}}\\
      \delta(&\hspace{-0.35cm} x\cdot z    \hspace{-0.35cm}&) = \delta(x) \land \delta(z)                        \\
      \delta(&\hspace{-0.35cm} x\vee  z    \hspace{-0.35cm}&) = \delta(x) \vee  \delta(z)                        \\
      \delta(&\hspace{-0.35cm} x\land z    \hspace{-0.35cm}&) = \delta(x) \land \delta(z)
    \end{eqnarray*}
  \end{multicols}
\end{definition}

Primarily, this gadget was designed to handle membership queries, for which purpose it has  found a number of applications~\cite{might2011parsing,stanford2021symbolic,varatalu2025re} in recent years:

\begin{theorem}[Recognition]\label{thm:recognition}
  For any regex \(e\) and \(\sigma: \Sigma^*\), \(\sigma \in \mathcal{L}(e) \Longleftrightarrow \varepsilon \in \mathcal{L}(\partial_\sigma e)\), where:

  \begin{equation}
    \partial_\sigma (e): E \rightarrow E = \begin{cases}e &\text{ if } \sigma = \varepsilon\\\partial_b(\partial_a e) &\text{ if } \sigma = a \cdot b, a \in \Sigma, b \in \Sigma^* \end{cases}
  \end{equation}
\end{theorem}

Variations on this basic procedure can also be used for functional parsing and regular expression tasks. Less well known, perhaps, is that Brzozowski's derivative can also be used to decode witnesses. We will first focus on the nonempty disjunctive fragment, and define this process in two steps:

\begin{theorem}[Generation]\label{thm:generation}
  For any nonempty $(\varepsilon, \land)$-free regex, \(e\), to witness $\sigma \in \mathcal{L}(e)$:

\begin{equation}
  \hspace{-1cm}\texttt{follow}(e): E \rightarrow 2^\Sigma = \begin{cases}
   \{e\} &\text{ if } e \in \Sigma \\
   \texttt{follow}(x) &\text{ if } e = x \cdot z\\
   \texttt{follow}(x)\cup\texttt{follow}(z) &\text{ if } e = x \lor z
  \end{cases}
\end{equation}

\begin{equation}\label{eq:choose}
  \texttt{choose}(e): E \rightarrow \Sigma^+ = \begin{cases}
   e &\text{ if } e \in \Sigma \\
   \big(s \leftsquigarrow \texttt{follow}(e)\big)\cdot \texttt{choose}(\partial_s e) &\text{ if } e = x \cdot z\\
   \texttt{choose}\big(e' \leftsquigarrow \{x, z\}\big) &\text{ if } e = x \lor z
  \end{cases}
\end{equation}
\end{theorem}

Here, we use the $\leftsquigarrow$ operator to denote probabilistic choice, however, any deterministic choice function will also suffice to generate a witness. Now we are equipped to handle conjunction.

Recall that every regular language is also context-free a fortiori. So, given an $(\varepsilon, \land)$-free regular expression, we can construct an equivalent CFG with productions $P(e)$ as follows:

\begin{equation}
P(e): E \rightarrow \big(V \rightarrow (\Sigma \mid V \mid V^2)\big) = \begin{cases}
 \{ S_e \rightarrow e \} & \text{if } e \in \Sigma \\
 P(x) \cup P(z) \cup \{ S_e \rightarrow S_x S_z \} & \text{if } e = x \cdot z \\
 P(x) \cup P(z) \cup \{ S_e \rightarrow S_x, S_e \rightarrow S_z \} & \text{if } e = x \lor z \\
\end{cases}
\end{equation}\vspace{0.2cm}

\noindent where the CFG is $G(e) = \langle V, \Sigma, P(e), S_e\rangle$ with $V$ being nonterminals in $P(e)$. Therefore, to intersect two regular languages, we can treat one of them as a CFL. Alternatively, we can take the intersection between some truly non-regular CFL (say, a programming language syntax) and a regular language.

\begin{theorem}[Bar-Hillel, 1961]
  For any CFG, $G = \langle V, \Sigma, P, S\rangle$, and nondeterministic finite automata (NFA), $A = \langle Q, \Sigma, \delta, q_\alpha, F\rangle$, there is a CFG, \(G_\cap=\langle V_\cap, \Sigma_\cap, P_\cap, S_\cap\rangle\) s.t. $\mathcal{L}(G_\cap) = \mathcal{L}(G)\cap\mathcal{L}(A)$.
\end{theorem}

\noindent Salomaa~\cite{salomaa1973formal} introduces a direct, but inefficient construction for the intersection grammar:

\begin{definition}[Salomaa, 1973]
  One could construct $G_\cap$ like so,

  \noindent\begin{prooftree}
      \hskip -0.5em
      \AxiomC{$q_\omega \in F\vphantom{\overset{a}{\rightarrow}}$}
      \RightLabel{$\mathcal{S}$}
      \UnaryInfC{$\big(S\rightarrow q_\alpha S q_\omega\big) \in P_\cap$}
      \DisplayProof
      \hskip 1em
      \AxiomC{$(w \rightarrow a) \in P$}
      \AxiomC{$(q\overset{a}{\rightarrow}r) \in \delta$}
      \RightLabel{$\uparrow$}
      \BinaryInfC{$\big(qwr\rightarrow a\big)\in P_\cap$}
      \DisplayProof
      \hskip 1em
      \AxiomC{$(w \rightarrow xz) \in P$}
      \AxiomC{$\vphantom{(}p,q,r \in Q\vphantom{\overset{a}{\rightarrow}}$}
      \RightLabel{$\Join$}
      \BinaryInfC{$\big(pwr\rightarrow (pxq)(qzr)\big) \in P_\cap$}
  \end{prooftree}
\end{definition}\vspace{0.2cm}

\noindent however, most synthetic productions in $P_\cap$ will be non-generating or unreachable. This method will construct a synthetic production for state pairs that are not even connected by any path, which is clearly excessive. In \S~\ref{sec:method}, we will present a far more efficient construction for the special case when the intersection is finite. But first, let us return to the broader question of syntax repair.

\subsection{Informal statement}

Assume there exists a transducer from Unicode tokens to grammatical tokens, $t: \Sigma_U^* \rightarrow \Sigma_G^*$. In the compiler nomenclature, $t$ is called a \textit{lexer} and would typically be regular under mild conditions. In this paper, we do not consider $t$ and strictly deal with languages over $\Sigma_G^*$, or simply $\Sigma^*$ for brevity.

Now suppose we have a syntax, $\ell \subset \Sigma^*$, containing every acceptable program. A syntax error is an unacceptable string, $\err\sigma \notin \ell$, that we wish to repair. We can model syntax repair as a language intersection between a context-free language (CFL) and a regular language. Henceforth, $\err\sigma$ will always and only be used to denote a syntactically invalid string whose intended language is known.

\begin{wrapfigure}{r}{0.4\textwidth}
\vspace{-0.3cm}
\resizebox{\linewidth}{!}{%
\begin{tikzpicture}[
  scale=4.5,
  >={Stealth[length=12pt, width=10pt]},
  dot/.style = {circle, inner sep=0pt, minimum size=4mm, fill, node contents={}},
  posterline/.style = {line width=3pt},
  labeltext/.style = {overlay, transform shape}
]
  % 1. Draw the underlying layers
  \draw[posterline] (-2.1,0) coordinate (a) circle (2.4) node[dot,label={[labeltext]above:$\err{\sigma}$},name=z0];

  \draw [posterline, dash pattern=on 15pt off 12pt] (a) circle (0.6);
  \draw [posterline, dash pattern=on 15pt off 12pt] (a) circle (1.2);
  \draw [posterline, dash pattern=on 15pt off 12pt] (a) circle (1.8);

  \draw[->, posterline] (-2.1,0) -- (-1.5, 0) node[midway,below,labeltext,xshift=-1pt,yshift=2pt]{$d_1$};
  \draw[->, posterline] (-1.5,0) -- (-0.9, 0) node[midway,below,labeltext,xshift=-1pt,yshift=2pt]{$d_2$};
  \draw[->, posterline] (-0.9,0) -- (-0.3, 0) node[midway,below,labeltext]{$d_3$};
  \draw[->, posterline] (-0.3,0) -- (0.3, 0) node[midway,below,labeltext]{$d_4$};

  % 2. Gray background for the intersection region
  \pgfmathsetseed{1234}
  \begin{scope}
    \clip (-2.1, 0) circle (2.4);
    \clip[decorate, decoration={random steps, segment length=15pt, amplitude=8pt}] (1.15, 0) circle (2.0);
    \fill[black!10] (-5,-5) rectangle (5,5);
  \end{scope}

  % 3. Draw the red dots strictly inside the mathematical boundary
  \begin{scope}
    \pgfmathsetseed{1234}
    \clip (-2.1, 0) circle (2.4);
    \clip[decorate, decoration={random steps, segment length=15pt, amplitude=8pt}] (1.15, 0) circle (2.0);

    \pgfmathsetseed{2673}
    \foreach \k in {1,...,400}{
      \pgfmathsetmacro{\x}{rnd*1.2 - 0.9}
      \pgfmathsetmacro{\y}{rnd*3.0 - 1.5}
      \fill[black] (\x, \y) circle[radius=0.017];
    }
  \end{scope}

  % 4. Redraw the boundaries
  \draw[posterline] (-2.1,0) circle (2.4);
  \pgfmathsetseed{1234}
  \draw[posterline, decorate, decoration={random steps, segment length=15pt, amplitude=8pt}] (1.15,0) circle (2.0);

  % 5. Draw external labels
  \node [above,labeltext] at (current bounding box.north -| a) {$\mathcal{L}\bigl(\text{L}(\err\sigma, d^* + 1)\bigr)$};
  \node [above,labeltext] at (1.15, 2.1) {$\mathcal{L}(G)$};
  \node [above,labeltext] at (-0.25, 1.8) {$\ell_\cap$};

  \path[use as bounding box] (-4.5,-2.4) rectangle (3.25,3.0);
\end{tikzpicture}%
}

\vspace{-0.3cm}
\caption{CFL intersection with the local edit region around a broken code snippet, where $d^*=3$ is the language edit distance (LED).}
\vspace{-0.3cm}
\end{wrapfigure}

Given a lexical representation of a broken computer program $\err\sigma$ and a grammar $G$, our goal is to find every valid string $\sigma$ consistent with the grammar $G$ and within a certain edit distance, $d$. Consider the language of nearby strings: if intersected with the language of grammatically valid programs, $\mathcal{L}(G)$, the result ($\ell_\cap$) will contain every possible repair within the given edit distance, a subset of which will be natural or statistically probable. If we can locate these repairs, then we can map them back into Unicode, adding placeholders for fresh names, numbers, and string literals, then finally apply an off-the-shelf code formatter to display them. Both the preprocessing and the cosmetic postprocessing steps are tangential to this work, in which we confine ourselves to a lexical alphabet.

\subsection{Formal statement}\label{sec:problem}

Let us now restate our informal description of the syntax repair problem in more formal terms.

\begin{definition}[Bounded Levenshtein-CFL reachability]\label{def:bcflr}
Given a CFL, $\ell$, and an invalid string, $\err{\sigma}: \bar\ell$, find every valid string reachable within $d$ edits of $\err{\sigma}$, i.e., letting $\Delta$ be the Levenshtein metric and $\mathcal{L}\big(L(\err\sigma, d)\big) = \{\sigma' \mid \Delta(\err{\sigma}, \sigma') \leq d\}$ be the Levenshtein $d$-ball, we seek to find $\ell_\cap = \mathcal{L}\big(L(\err\sigma, d)\big) \cap \ell$.
\end{definition}

As the intersection language, $\ell_\cap$, typically contains a large number of possible repairs, we want a procedure that surfaces both natural and valid repairs over unnatural but valid repairs:

\begin{definition}[Ranked repair]\label{def:ranked-repair}
Given a finite language $\ell_\cap = \mathcal{L}\big(L(\err\sigma, d)\big) \cap \ell$ and a probabilistic language model $\text{P}_\theta: \Sigma^* \rightarrow [0, 1] \subset \mathbb{R}$, find the top-$k$ maximum probability repairs. That is,
\begin{equation}
R(\ell_\cap, P_\theta): 2^{\Sigma^*} \times (\Sigma^* \rightarrow \mathbb{R}) \rightarrow (\Sigma^*)^{\leq k} = \argmax_{\bm{\sigma} \subseteq \ell_\cap, |\bm{\sigma}| \leq k} \sum_{\sigma \in \bm{\sigma}}\text{P}_\theta(\sigma)
\end{equation}
\end{definition}

A popular approach to ranked repair involves learning a distribution over strings, however, this is highly sample-inefficient and generalizes poorly to new languages. Approximating a distribution over $\Sigma^*$ forces the model to jointly learn syntax and stylometry. Furthermore, even with an extremely efficient approximate sampler for $\sigma \leftsquigarrow \ell_\cap$, due to the size of the languages involved, it would be intractable to sample either $\ell$ or $\mathcal{L}\big(L(\err\sigma, d)\big)$, reject duplicates, then reject unreachable or invalid edits, and completely out of the question to sample $\sigma \leftsquigarrow \Sigma^*$, as do most neural language models.

As we will demonstrate, the ranked repair problem can be factorized into three subproblems: (1) exact representation, (2) decoding and (3) reranking. Instead of working with strings, we will explicitly construct a grammar which soundly and completely generates $\ell \cap \mathcal{L}\big(L(\err\sigma, d)\big)$, then decode repairs from its language. By ensuring decoding is sufficiently precise and extensive, ensuring the retrieved set contains the true repair can be achieved with a much simpler, syntax-oblivious model. Finally, we will train a language model to rerank the repair candidates and take the top-$k$ results.
\ifSubfilesClassLoaded{\par\clearpage\bibliography{../bib/acmart}}{}
\end{document}
